\scanpage{9}
% slanted
\margin{\deleted{\unsure{que}} Précisons que tout point de l'icosaèdre est centre d'une \deleted{\unsure{déterminée}} unique calotte d'icosaèdre contenue dans $I$ ; la calotte de centre $a$ est formée de $a$ et de tous les pts \uline{reliés} à $a$ par une arrête. \deleted{\unsure{Elle est}} \deleted{\unsure{unique}}}
\margin{Données équivalentes :
$\left\{\begin{array}{l}
\text{icosaèdre + un des sommets}\\
\text{icosaèdre + calotte d'icosaèdre contenue dedans}\\
\text{calotte d'icosaèdre}\\
\text{une couple }(u,v)\text{ modulo rotation sur }v\\
\quad\text{dans }u^\perp\text{ d'angle multiple de }\theta
\end{array}\right.$
[point du fibré principal sur $S$ de groupe $\mathbb T$ \unsure{de|,} puissance cinquième du fibré fourni \unsure{des|du} \unsure{$(u,v)$|$u,v$}, \unsure{soit|\ldots}}

de $u$ dans \deleted{\unsure{G}} $\Gamma'$ est aussi le sous-groupe \deleted{de} $G'$ \deleted{formé des} des $g\in G'$ qui invarient la calotte $(u,(x_i))$, lequel, nous avons vu, est le groupe $\mu_5$ qui opère de façon simplement transitive sur l'ens des $x_i$, OK.

Prouvons enfin l'assertion de \inserted{simple} transitivité de $\Gamma$ sur l'ens. des \deleted{faces} triples $(x,y,z)$ qui sont de face. Procédant comme dessus, on peut supposer, utilisant ce qui est déjà démontré, que $x=x'=u,y=y'=x_0,z=x_1$. Mais il n'y a que deux \deleted{triangles} faces \deleted{\unsure{en}} admettant $(u,x_0)$ comme arrête, leurs autre sommet est resp. $x_1$ et $x_{-1}$. Si c'est $x_{-1}$, la symétrie p.r. au plan $u,v$ (qui évidemment invarie l'icosaèdre $I$) transforme $x_{-1}$ en $x_1$ en laissant $u,x_0$. Il reste donc : \deleted{vérifier} noter que $u,x_0,x_1$ \deleted{\unsure{x_{-1}}} étant linéairement indépendants une transformation linéaire qui fixe ces trois points est l'identité.

\uline{Corollaire} $\Gamma'$ est d'ordre 60, $\Gamma$ d'ordre 120.

Le premier point provient du fait que l'icosaèdre a 30 arrêtes, donc $2.30=60$ arrêtes orientées. Comme $\Gamma\simeq\Gamma'\times\mu_2$, on trouve que $\Gamma$ a 120 éléments. On peut aussi noter qu'il y a 20 faces, chacune ayant \deleted{\unsure{deux|trois}} six ordres totaux sur l'ens de ses sommets, ce qui fait $6\times20=120$ triples ordonnés faisant face.

\uline{Remarque} Les sommets et arrêtes de l'icosaèdre forment (avec leur relation d'incidence) un \uline{graphe}, et $\Gamma'$ opère par automorphismes de ce graphe. Montrons que \uline{tout} automorphisme \inserted{$g$} du graphe est dans $\Gamma'$. On est ramené à \deleted{\unsure{en}} prouver \deleted{\unsure{pour}} qu'un autom. du graphe qui fixe $u,x_0,x_1$ \deleted{\unsure{g}} est l'identité. Prouvons qu'il est l'identité sur tous ces points $x_i$ (\inserted{i.e. tous les pts de $I$,} qui peuvent être reliés à \unsure{l'un|un} des sommets -- $u$ en l'occurrence -- par une arrête) : en effet \deleted{\unsure{on}} pour $x_0,x_1$ il n'y a rien à prouver, or $x_2$ ($x_{-1}$) est l'unique point de $I$ \inserted{distinct de $x_0$ (resp. $x_1$)} qui est relié \deleted{par} à la fois à $u$ et $x_1$ ($u$ et $x_0$), donc il est fixé, donc aussi $x_3$ (p. ex. comme
% archive sheet number: 8
